\section{Semi Parametric KDE}
\begin{frame}{\hspace{3cm}}
\centering

\vspace{1.5cm}

\begin{beamercolorbox}[sep=20pt,center,rounded=true,shadow=true]{title}
    {\fontsize{20}{24}\selectfont \textbf{Semi-Parametric KDE}}
\end{beamercolorbox}

\vspace{0.8cm}


\vfill

\end{frame}
%------------------------------------------------
\begin{frame}{Semi-Parametric KDE}
Suppose that 
\[
X := (X_1, \ldots, X_d), \quad i = 1, \ldots, n,
\]
is a random sample of size $n$ from a $d$-dimensional density function, and let 
\[
\mathbf{X}_i = (X_{i1}, \ldots, X_{id})
\]
be an observation.

\textit{Azzalini and Menardi (2014)} performed the density estimation step of their clustering algorithm using a kernel method based on a product kernel estimator of the form:
\begin{equation}
\hat{f}(\mathbf{x}) = \frac{1}{n} \sum_{i=1}^{n} \prod_{j=1}^{d} \frac{1}{h_{ij}} \, \kappa\!\left( \frac{x_j - X_{ij}}{h_{ij}} \right),
\end{equation}
where
\[
\mathbf{x} = (x_1, \ldots, x_d),
\]
$\kappa(\cdot)$ is a symmetric kernel (e.g., normal or $t_7$ as used in \texttt{pdfCluster}), and $h_{ij}$ denotes the bandwidth corresponding to $X_{ij}$ for $i = 1, \ldots, n$ and $j = 1, \ldots, d$.
    
\end{frame}
\begin{frame}{Proposed Skew-Symmetric Kernel Estimator}
Here, we intend to employ a product kernel of the following form:
\begin{equation}
\hat{f}(\mathbf{x}) = \frac{1}{n \, h_1 \cdots h_d} \sum_{i=1}^{n} \prod_{j=1}^{d} 
\kappa_j \!\left( \frac{x_j - X_{ij}}{h_j} \right),
\end{equation}
where $\kappa_j(\cdot)$ can be either a skew-normal density or a skew-$t$ density.

More precisely, in the skew-normal case, we have:
\begin{equation}
\kappa_j(z) = 2 \, \phi(z)\, \Phi(\lambda_j z), \quad z \in \mathbb{R},
\end{equation}
where $\lambda_j$, $j = 1, \ldots, d$, is the slant parameter, and $\phi(\cdot)$ and $\Phi(\cdot)$ denote the density and cumulative distribution function of the standard normal distribution, respectively.
    
\end{frame}
\begin{frame}{Proposed Skew-Symmetric Kernel Estimator(Contd.)}
Our second proposal for $\kappa_j(\cdot)$ is the standard skew-$t$ density given by:
\begin{equation}
\kappa_j(z) = 2 \, t(z; \nu_j)\, 
T\!\left( \lambda_j z \sqrt{\frac{\nu_j + 1}{\nu_j + z^2}} \, ; \, \nu_j + 1 \right), 
\quad z \in \mathbb{R},
\end{equation}
where $\lambda_j$ is the slant parameter and $\nu_j$, $j = 1, \ldots, d$, denotes the degrees of freedom.

Here, $t(\cdot \, ; \, \nu)$ and $T(\cdot \, ; \, \nu)$ represent the density and cumulative distribution function of the classical Student's $t$ distribution with $\nu$ degrees of freedom, respectively.
\end{frame}
\begin{frame}{Our Algorithm}
Let 
\[
\mathbf{z}_j := (Z_{1j}, \ldots, Z_{nj}).
\]

\begin{itemize}
    \item Choose the desired kernel between (3) and (4).

    \item For the skew-normal (SN) kernel, plug in $\lambda_j$ using its MPLE, assuming $\mathbf{z}_j$ is a random sample from (5). \textit{(A point on MPLE)}

    \item For the skew-$t$ (ST) kernel, proceed as follows, assuming $\mathbf{z}_j$ is a random sample from (6):
    \begin{itemize}
        \item[(i)] For high dimensions: approximate $(\lambda_j, \nu_j)$ with quick preliminary estimates using \texttt{st.prelimFit}.
        
        \item[(ii)] For low dimensions: plug in $(\lambda_j, \nu_j)$ using their MPLEs via the routines \texttt{st.prelimFit} and \texttt{st.mple}.
    \end{itemize}

    \item Compute $\hat{f}(\mathbf{x})$ in (4) with $\kappa_j(\cdot)$ whose parameters are now approximated.

    \item The only remaining ingredient is to determine an optimal bandwidth, which will be discussed in the next section.
\end{itemize}
    
\end{frame}
\begin{frame}{On finding optimum bandwidth}

Suppose that $f(\cdot)$ is a univariate density function to be estimated using (2), considering $d = 1$.  
Assume further that $\kappa(\cdot)$ is a one-dimensional asymmetric kernel density satisfying the following regularity assumptions:

\vspace{0.4cm}

\begin{block}{Regularity Assumptions}
\begin{itemize}
    \item[(A$_1$)] $z \in \mathbb{R}, \quad \kappa(z) \geq 0$
    
    \item[(A$_2$)] $\displaystyle \int \kappa(z)\,dz = 1$
    
    \item[(A$_3$)] $\displaystyle \int z\,\kappa(z)\,dz = \mu_1(\kappa) \neq 0$
\end{itemize}
\end{block}

\end{frame}
\begin{frame}{Optimal Bandwidth $h_{\text{opt}}$}

Now, if we consider $\kappa(\cdot)$ to be the one given by (3), then $h_{\text{opt}}$ is obtained as a real value satisfying:

\vspace{0.3cm}

\begin{block}{Equation}
\[
-a + 2b\,h_{\text{opt}}^{3}(\lambda) + c\,h_{\text{opt}}^{5}(\lambda) = 0
\tag{5}
\]
\end{block}
where,


\[
a := \frac{1}{n}\left(2\pi^\frac{-3}{2}\arctan\!\left(\sqrt{1+\lambda^{2}}\right)\right)
\]

\[
b := \frac{\lambda^{2}}{2\pi^{3/2}\,\sigma^{3}(1+\lambda^{2})}
\]

\[
c := \frac{3}{8\pi^{1/2}\,\sigma^{5}}
\tag{6}
\]


\end{frame}
\begin{frame}
\begin{figure}
        \centering
        \includegraphics[width= 0.63\linewidth]{Screenshot 2026-04-14 085621.png}
        \label{fig:placeholder}
    \end{figure}    
\end{frame}
\begin{frame}{Multi-dimensional Case}

Assume $f(\mathbf{x})$ is a $d$-variate density. Then, the kernel density estimator in its most general form is defined as \\
\textit{[Wand and Jones (1995)]}:

\vspace{0.4cm}

\begin{block}{Kernel Density Estimator}
\[
\hat{f}(\mathbf{x}) = \frac{1}{n} |\mathbf{H}|^{-1/2} \sum_{i=1}^{n} 
K\!\left(\mathbf{H}^{-1/2}(\mathbf{x} - \mathbf{X}_i)\right)
\tag{7}
\]
\end{block}

\vspace{0.3cm}

\begin{block}{Notation}
\begin{itemize}
    \item $\mathbf{x} = (x_1, \dots, x_d)$
    \item $\mathbf{X}i = (X{i1}, \dots, X_{id}), \quad i = 1, \dots, n$
    \item $\mathbf{H}$: bandwidth matrix (symmetric positive definite $d \times d$ matrix)
    \item $K(\cdot)$: $d$-variate kernel density
\end{itemize}
\end{block}

\end{frame}
\begin{frame}{Theorem: AMISE in Multi-dimensional Case}

\begin{block}{Theorem}
Assume $h_j = h$, for $j = 1, \dots, d$. Then, the AMISE of $\hat{f}(\mathbf{x})$ has the form:
\[
\mathrm{AMISE}(\hat{f}) = \frac{1}{n h^{d}} \int t^{1/3}(x)\,dx \;+\; A\left(Bh^{2} + Ch^{4}\right)
\tag{8}
\]

\end{block}
where,
\vspace{-3 mm}
\[
B := 8 \sum_{j=1}^{d} \mu_{1}^{2}(\kappa_j)\,\sigma_{j}^{-2}
\]
\vspace{-3 mm}
\[
A := (2\sqrt{\pi})^{-d} \prod_{j=1}^{d} \sigma_{j}^{-1}
\]
\vspace{-3 mm}
\[
C := 3 \sum_{j=1}^{d} \mu_{2}^{2}(\kappa_j)\,\sigma_{j}^{-4}
+ 2 \sum_{j<l} \left[
\mu_{2}(\kappa_j)\mu_{2}(\kappa_l) 
+ \mu_{1}^{2}(\kappa_j)\mu_{1}^{2}(\kappa_l)
\right]\sigma_{j}^{-2}\sigma_{l}^{-2}
\]


\end{frame}
\begin{frame}
Minimizing (8) gives the optimum bandwidth as the solution of the following non-linear equation
\begin{equation}
- \frac{d}{n} \prod_{j=1}^{d} \int \kappa_j^2(t)\, dt 
+ 2AB h_{\mathrm{opt}}^{d+2} 
+ 4AC h_{\mathrm{opt}}^{d+4} = 0,
\tag{9}
\end{equation}

where constants $A$, $B$ and $C$ are given by (9). The unknown parameters $\sigma_j$ can be substituted by $s_j$, the sample standard deviation of the $j$th variable, and the $\lambda_j$ can be approximated by its MPLE.

It can be easily shown that by setting $d = 1$ in (9), it simplifies to its univariate version given by (5).
    
\end{frame}
\begin{frame}
\begin{figure}
    \centering
    \includegraphics[width=0.75\linewidth]{Screenshot 2026-04-14 100602.png}
    \label{fig:placeholder}
\end{figure}
    
\end{frame}
\begin{frame}
\begin{figure}
    \centering
    \includegraphics[width=0.75\linewidth]{Screenshot 2026-04-14 103509.png}
    \label{fig:placeholder}
\end{figure}
    
\end{frame}


%-------------------------------
